Let $\mathbf{X}_1,\ldots,\mathbf{X}_n$ be independent copies of a random vector
$\mathbf{X}\in\mathcal X\subseteq\mathbb R^p$ with distribution $P_0$.  We do not
initially require $P_0$ to belong to a finite-mixture model.  This distinction
allows the classification likelihood to be treated as a possibly
misspecified criterion and makes its population target explicit.

Consider $k\geq2$ positive component scores
\begin{equation}
 q_j(\mathbf{x};\boldsymbol{\theta})=\log\pi_j(\boldsymbol{\theta})+\log f_j(\mathbf{x};\boldsymbol{\alpha}_j(\boldsymbol{\theta})),
 \qquad j=1,\ldots,k,
 \label{eq:component-score}
\end{equation}
where $f_j(\cdot;\boldsymbol{\alpha}_j)$ may belong to different parametric families.  The
component parameter $\boldsymbol{\alpha}_j$ has dimension $d_j$.  The full parameter
$\boldsymbol{\theta}$ is expressed in a local unconstrained coordinate system and takes
values in $\Theta\subset\mathbb R^d$, where
\begin{equation}
 d=(k-1)+\sum_{j=1}^k d_j.
 \label{eq:parameter-dimension}
\end{equation}
For example, the first $k-1$
coordinates may be multinomial-logit coordinates for the mixing weights,
while covariance matrices may be represented by unconstrained Cholesky
coordinates.  All derivatives below are with respect to this fixed coordinate
system.

Define
\begin{equation}
 m_{\boldsymbol{\theta}}(\mathbf{x})=\max_{1\leq j\leq k}q_j(\mathbf{x};\boldsymbol{\theta}),\qquad
 M_n(\boldsymbol{\theta})=\mathbb P_n m_{\boldsymbol{\theta}},\qquad
 M(\boldsymbol{\theta})=P_0m_{\boldsymbol{\theta}},
 \label{eq:cml-criteria}
\end{equation}
where $\mathbb P_n=n^{-1}\sum_{i=1}^n\delta_{\mathbf{X}_i}$.  The population
classification maximum-likelihood target is
\begin{equation}
 \boldsymbol{\theta}^*=\arg\max_{\boldsymbol{\theta}\in\Theta}M(\boldsymbol{\theta}).
 \label{eq:cml-target}
\end{equation}
The equality in \eqref{eq:cml-target} defines a statistical functional of
$P_0$; it does not assert that $\boldsymbol{\theta}^*$ is the parameter generating $P_0$.
Even under a correctly specified ordinary mixture model, the CML target and
the ordinary mixture-likelihood parameter need not coincide.

The theoretical estimator is any measurable approximate maximizer satisfying
\begin{equation}
 M_n(\widehat{\boldsymbol{\theta}}_n)
 \geq \sup_{\boldsymbol{\theta}\in\Theta}M_n(\boldsymbol{\theta})-r_n,
 \qquad r_n=o_p(1).
 \label{eq:approximate-maximizer}
\end{equation}
Exact empirical maximizers correspond to $r_n=0$.  Stronger control of $r_n$
will be imposed when deriving a first-order distribution.  Condition
\eqref{eq:approximate-maximizer} is a property to be verified for an
algorithmic output; calling an arbitrary stationary point a CML estimator does
not make the condition true.

For $j\ne l$, write
\begin{equation}
 g_{jl}(\mathbf{x};\boldsymbol{\theta})=q_j(\mathbf{x};\boldsymbol{\theta})-q_l(\mathbf{x};\boldsymbol{\theta}).
 \label{eq:pairwise-contrast}
\end{equation}
Let $a_{\boldsymbol{\theta}}(\mathbf{x})$ be the smallest index in
$\arg\max_jq_j(\mathbf{x};\boldsymbol{\theta})$.  This deterministic convention makes
$a_{\boldsymbol{\theta}}$ measurable and defines the disjoint classification cells
\begin{equation}
 \mathcal R_j(\boldsymbol{\theta})=\{\mathbf{x}:a_{\boldsymbol{\theta}}(\mathbf{x})=j\},\qquad j=1,\ldots,k.
 \label{eq:classification-cells}
\end{equation}
The convention has no effect on population quantities whenever ties have
$P_0$-probability zero.

\subsection{Parameter identifiability}

When two components use the same family, simultaneous permutations of their
weights and component parameters leave $m_{\boldsymbol{\theta}}$ unchanged.  We therefore
work on an identifiable representative set $\Theta$.  This can be obtained by
a scientifically meaningful component ordering, a fixed family assignment in
a heterogeneous model, or a canonical ordering of a component functional.
The chosen ordering must hold in a neighbourhood of $\boldsymbol{\theta}^*$; ordering
components after estimation is not a substitute for local identifiability.

The ordinary mixture density
$\sum_j\pi_jf_j(\cdot;\boldsymbol{\alpha}_j)$ is used for data generation in some
experiments, but it is not used to replace the unknown $P_0$ in the general
theory.  In particular, the feasible boundary estimator developed later is
based on observations from $P_0$ and does not assume that the mixture evaluated
at $\boldsymbol{\theta}^*$ equals the data-generating density.

\begin{assumption}[Basic criterion conditions]
\label{ass:basic-criterion}
The following conditions hold.
\begin{enumerate}
\item $\Theta$ is a compact, identifiable subset of $\mathbb R^d$.
 \item For every $j$, $(\mathbf{x},\boldsymbol{\theta})\mapsto q_j(\mathbf{x};\boldsymbol{\theta})$ is jointly
 measurable, and $\boldsymbol{\theta}\mapsto q_j(\mathbf{x};\boldsymbol{\theta})$ is continuous for
 $P_0$-almost every $\mathbf{x}$.
 \item There is a measurable function $F$ with $P_0F<\infty$ such that
 $\sup_{\boldsymbol{\theta}\in\Theta}\max_j|q_j(\mathbf{X};\boldsymbol{\theta})|\leq F(\mathbf{X})$ almost surely.
 \item $M$ has a unique, well-separated maximizer $\boldsymbol{\theta}^*$: for every
 $\varepsilon>0$,
 \[
  M(\boldsymbol{\theta}^*)-
  \sup_{\boldsymbol{\theta}\in\Theta:\,\|\boldsymbol{\theta}-\boldsymbol{\theta}^*\|\geq\varepsilon}M(\boldsymbol{\theta})>0.
 \]
\end{enumerate}
\end{assumption}

Compactness is used only for the global consistency argument.  Later local
results require instead that $\boldsymbol{\theta}^*$ be an interior point of an open
coordinate neighbourhood and that the empirical procedure enters that
neighbourhood with probability tending to one.
